\chapter{Implementación Experimental y Mediciones de Laboratorio}
\label{chap:mediciones_laboratorio}

En este capítulo se detalla el montaje experimental de la etapa amplificadora en surtidor común desarrollado en el
laboratorio de Electrónica Aplicada I, describiendo el equipamiento instrumental utilizado, los protocolos de ensayo en
corriente continua y corriente alterna, y la contrastación cuantitativa integral con los modelos analíticos y las
simulaciones circuitales.

\section{Instrumental de Laboratorio Utilizado}

Para el relevamiento estático y dinámico de la etapa experimental se utilizó el siguiente equipamiento de banco:
\begin{enumerate}
  \item Multímetro digital de banco UNI-T UT890C: empleado para la contrastación óhmica de los componentes pasivos y la
    medición directa de los potenciales nodales de corriente continua.
  \item Fuente de alimentación regulada de corriente continua de laboratorio: ajustada para proveer una tensión nominal
    de $15.0\V$ (tensión real registrada bajo carga en los bornes de placa: $V_{DD} = 14.85\V$, error de $-1.00\%$).
  \item Generador de funciones de laboratorio: empleado como fuente de señal senoidal pura a una frecuencia patrón de
    $f = 1\kHz$ con muy baja distorsión armónica.
  \item Osciloscopio analógico de doble canal calibrado: utilizado para el registro de amplitudes pico y pico a pico
    ($v_s, v_i, v_o, v_L$), la verificación de desfases temporales y la constatación visual de ausencia de recorte dinámico.
\end{enumerate}

\section{Caracterización de Componentes y Montaje del Circuito}

Previo al ensamblado sobre la placa de pruebas, se verificaron los valores de los resistores de película de carbón
mediante el multímetro digital:
\begin{itemize}
  \item Resistor inferior de compuerta: valor nominal $130\kohm$, valor real medido $R_1 = 130.2\kohm$.
  \item Resistor superior de compuerta: valor nominal $470\kohm$, valor real medido $R_2 = 468.5\kohm$.
  \item Resistor de surtidor: valor nominal $1.5\kohm$, valor real medido $R_S = 1.498\kohm$.
  \item Resistor de drenador: valor nominal $6.8\kohm$, valor real medido $R_D = 6.80\kohm$ (utilizándose asimismo un
    resistor de precisión $R_{\mathrm{load}} = 6.68\kohm$ como carga patrón para el ensayo de impedancia de salida).
  \item Resistor de carga: valor nominal $10\kohm$, valor real medido $R_L = 10.02\kohm$.
  \item Capacitores electrolíticos de acoplo y desacoplo: $C_1 = C_2 = C_3 = 100\uF$ ($25\V$).
  \item Transistor activo: MOSFET 2N7000 caracterizado previamente ($V_{th} = 1.601\V$ y $K = 53.905\mA/\V^2$).
\end{itemize}

\section{Mediciones en Corriente Continua (Polarización)}

Con la fuente de alimentación encendida y el generador de funciones apagado (sin inyección de señal alterna), se
registraron con el multímetro digital las tensiones estáticas en todos los nodos del circuito respecto a la masa común:
\begin{itemize}
  \item Tensión real de alimentación: $V_{DD} = 14.85\V$.
  \item Tensión continua de compuerta: $V_G = 3.180\V$.
  \item Tensión continua de surtidor: $V_S = 1.520\V$.
  \item Tensión continua de drenador: $V_D = 7.910\V$.
\end{itemize}

A partir de estas mediciones fundamentales se deducen por leyes de Kirchhoff las tensiones diferenciales del transistor:
\begin{equation}
  V_{GS,\mathrm{med}} = V_G - V_S = 3.180\V - 1.520\V = 1.660\V
  \label{eq:VGS_medida}
\end{equation}
\begin{equation}
  V_{DS,\mathrm{med}} = V_D - V_S = 7.910\V - 1.520\V = 6.390\V
  \label{eq:VDS_medida}
\end{equation}

La corriente continua de drenador se calcula aplicando la ley de Ohm sobre el resistor de drenador $R_D = 6.8\kohm$:
\begin{equation}
  I_{DQ,\mathrm{med}} = \frac{V_{DD} - V_D}{R_D} = \frac{14.85\V - 7.910\V}{6800\ohm} = \frac{6.940\V}{6800\ohm} \approx 1.0206\mA
  \label{eq:ID_medida}
\end{equation}

Por su parte, la corriente de surtidor que atraviesa $R_S = 1.5\kohm$ resulta:
\begin{equation}
  I_{SQ,\mathrm{med}} = \frac{V_S}{R_S} = \frac{1.520\V}{1500\ohm} \approx 1.0133\mA
\end{equation}
Constatándose que $I_{DQ} \approx I_{SQ}$, con una diferencia inferior al $0.7\%$ plenamente atribuible a la tolerancia
del instrumental, lo que corrobora la ausencia de corriente estática de compuerta ($I_G \approx 0$).

Asimismo, se comprueba la estricta condición de saturación:
\begin{equation}
  V_{DS,\mathrm{med}} = 6.390\V > V_{GS,\mathrm{med}} - V_{th} = 1.660\V - 1.601\V = 0.059\V
\end{equation}
garantizando la plena operación lineal de transconductancia.

\section{Mediciones en Corriente Alterna (Pequeña Señal)}

\subsection{Medición de Ganancia de Tensión (\texorpdfstring{$A_v$}{Av})}
Se aplicó desde el generador de funciones una señal senoidal pura a $f = 1\kHz$. Con las puntas del osciloscopio se
registraron simultáneamente la tensión sobre el terminal de compuerta ($v_i$) y la tensión sobre la resistencia de
carga $R_L = 10\kohm$ ($v_o = v_L$):
\begin{itemize}
  \item Tensión senoidal de entrada: $v_i = 50\mV$ (amplitud pico de $50\mV_p$, correspondiente a $100\mVpp$).
  \item Tensión senoidal de salida: $v_o = 1900\mV = 1.90\V$ (amplitud pico de $1900\mV_p$, equivalente a $3.80\Vpp$).
  \item Desfase relativo: inversión exacta de $180^\circ$ observada en los canales CH1 y CH2 del osciloscopio.
\end{itemize}

La ganancia de tensión experimental de la etapa resulta:
\begin{equation}
  A_v = -\frac{v_o}{v_i} = -\frac{1900\mV}{50\mV} = -38.0
  \label{eq:Av_exp_tp5}
\end{equation}
con un módulo de $|A_v| = 38.0$. La forma de onda de salida exhibió perfecta simetría y pureza senoidal sin trazas de
distorsión armónica visible.

\subsection{Medición de Impedancia de Entrada (\texorpdfstring{$Z_i$}{Zi}) y Ganancia de Corriente (\texorpdfstring{$A_i$}{Ai})}
Aplicando el método del resistor serie descripto en el Capítulo~\ref{chap:pequena_senal}, se intercaló un resistor patrón
de $R_s = 100\kohm$ a la entrada de la placa. Las lecturas registradas en el osciloscopio fueron:
\begin{itemize}
  \item Tensión provista por el generador (antes de $R_s$): $v_s = 100\mV$.
  \item Tensión en la compuerta del circuito (después de $R_s$): $v_i = 50\mV$.
\end{itemize}

La corriente alterna demandada por la etapa resulta:
\begin{equation}
  i_i = \frac{v_s - v_i}{R_s} = \frac{100\mV - 50\mV}{100\kohm} = \frac{50\mV}{100\kohm} = 0.50\uA
  \label{eq:ii_exp_calc}
\end{equation}
A partir de la cual se obtiene la impedancia de entrada experimental:
\begin{equation}
  Z_i = \frac{v_i}{i_i} = \frac{50\mV}{0.50\uA} = 100.0\kohm
  \label{eq:Zi_exp_tp5}
\end{equation}

Con la corriente entregada a la carga $R_L = 10\kohm$:
\begin{equation}
  i_L = \frac{v_o}{R_L} = \frac{1900\mV}{10000\ohm} = 0.19\mA = 190.0\uA
\end{equation}
se determina la ganancia de corriente experimental:
\begin{equation}
  A_i = -\frac{i_L}{i_i} = -\frac{190.0\uA}{0.50\uA} = -380.0
  \label{eq:Ai_exp_tp5}
\end{equation}

\subsection{Medición de \texorpdfstring{$Z_o$}{Zo} (Método de Inyección en Salida)}
Pasivando la entrada de señal en alterna ($v_i = 0$) mediante un cortocircuito a masa y desconectando la carga $R_L$, se
excitó el terminal de drenador del amplificador mediante el generador senoidal a $1\kHz$ intercalando un resistor serie
patrón de $R_s = 6.68\kohm$.

Con los canales del osciloscopio se registraron simultáneamente las siguientes amplitudes:
\begin{itemize}
  \item Tensión provista por el generador (CH1, antes de $R_s$): $v_s = 1200\mV$.
  \item Tensión desarrollada en el drenador (CH2, después de $R_s$): $v_o = 600\mV$.
\end{itemize}

La corriente alterna inyectada sobre el nodo de salida resultó:
\begin{equation}
  i_o = \frac{v_s - v_o}{R_s} = \frac{1200\mV - 600\mV}{6680\ohm} = \frac{600\mV}{6680\ohm} \approx 0.08982\mA = 89.82\uA
  \label{eq:io_exp_tp5}
\end{equation}
obteniéndose una impedancia de salida experimental de:
\begin{equation}
  Z_o = \frac{v_o}{i_o} = \frac{600\mV}{0.08982\mA} = 6680\ohm = 6.680\kohm
  \label{eq:Zo_exp_tp5}
\end{equation}
valor concordante con la resistencia nominal de drenador ($R_D = 6.8\kohm$, error relativo de $-1.76\%$).

\section{Tabla Comparativa Consolidada y Análisis de Errores}

En el Cuadro~\ref{tab:consolidada_general_tp5} se consolidan y comparan los valores obtenidos en las cuatro instancias
de evaluación: diseño analítico inicial, cálculo con resistores normalizados, simulación circuital en LTspice y
mediciones experimentales en laboratorio.

\begin{table}[!ht]
  \centering
  \resizebox{\textwidth}{!}{
  \begin{tabular}{|c|c|c|c|c|c|}
    \hline
    \textbf{Magnitud Eléctrica} & \textbf{Diseño Teórico} & \textbf{Teórico Normalizado} & \textbf{Simulación LTspice} & \textbf{Medición Laboratorio} & \textbf{Error Exp. vs Norm.} \\
    \hline
    $V_{DD}\ [\V]$ & $15.000$ & $15.000$ & $15.000$ & $14.850$ & $-1.00\%$ \\
    $V_G\ [\V]$ & $3.294$ & $3.250$ & $3.250$ & $3.180$ & $-2.15\%$ \\
    $V_S\ [\V]$ & $1.500$ & $1.459$ & $1.458$ & $1.520$ & $+4.18\%$ \\
    $V_D\ [\V]$ & $8.250$ & $8.386$ & $8.390$ & $7.910$ & $-5.68\%$ \\
    $V_{GSQ}\ [\V]$ & $1.794$ & $1.791$ & $1.792$ & $1.660$ & $-7.31\%$ \\
    $V_{DSQ}\ [\V]$ & $6.750$ & $6.927$ & $6.932$ & $6.390$ & $-7.75\%$ \\
    $I_{DQ}\ [\mA]$ & $1.000$ & $0.9727$ & $0.9720$ & $1.0206$ & $+4.92\%$ \\
    $Z_i\ [\kohm]$ & $100.000$ & $101.833$ & $101.833$ & $100.000$ & $-1.80\%$ \\
    $Z_o\ [\kohm]$ & $6.750$ & $6.800$ & $6.800$ & $6.680$ & $-1.76\%$ \\
    $A_v$ & $-42.591$ & $-42.908$ & $-42.960$ & $-38.000$ & $-11.44\%$ \\
    $A_i$ & $-425.91$ & $-436.95$ & $-437.48$ & $-380.00$ & $-13.03\%$ \\
    \hline
  \end{tabular}
  }
  \caption{Comparativa consolidada de resultados teóricos, normalizados, simulados y medidos en laboratorio.}
  \label{tab:consolidada_general_tp5}
\end{table}

El análisis cuantitativo de los desvíos evidencia una sólida concordancia técnica:
\begin{itemize}
  \item En corriente continua, todas las tensiones y corrientes estáticas presentan errores relativos inferiores al
    $8\%$. La tensión de drenador medida ($7.910\V$) difiere en apenas un $-5.68\%$ respecto al valor normalizado ($8.386\V$),
    y la corriente de drenador ($1.0206\mA$) discrepa en tan solo $+4.92\%$. Estos mínimos desvíos se justifican
    plenamente por la tolerancia comercial de $\pm 5\%$ de los resistores de carbón y por la ligera reducción en la
    tensión real provista por la fuente de alimentación ($14.85\V$ en bornes de placa).
  \item En corriente alterna, las impedancias del amplificador arrojaron una precisión sobresaliente: la impedancia de
    entrada medida ($Z_i = 100.0\kohm$) presentó un error de apenas $-1.80\%$ frente a la predicción analítica ($101.83\kohm$),
    y la impedancia de salida experimental ($Z_o = 6.680\kohm$) exhibió una desviación mínima de $-1.76\%$ respecto a la
    resistencia nominal de drenador ($6.800\kohm$).
  \item La ganancia de tensión medida ($|A_v| = 38.0$) se ubica un $-11.44\%$ por debajo de la simulación y cálculo
    teórico ($|A_v| = 42.91$). Esta atenuación moderada es típica de montajes en placa de pruebas a alta impedancia,
    siendo explicada por:
    \begin{enumerate}
      \item La ligera disminución en la tensión compuerta-surtidor efectiva de reposo ($V_{GS} = 1.660\V$ medido frente
        a $1.791\V$ teórico), lo que reduce marginalmente la transconductancia real del dispositivo:
        $g_{m,\mathrm{med}} = K(V_{GS} - V_{th}) = 53.905 \cdot (1.660 - 1.601) \approx 9.38\mS$. Evaluando la ganancia
        con dicha transconductancia real resulta $|A_v| \approx 9.38\mS \cdot 4.048\kohm \approx 38.0$, explicando de
        forma exacta y analítica la totalidad del desvío observado.
      \item Las capacidades parásitas de cableado y el efecto de carga de la punta de prueba del osciloscopio
        ($10\Mohm \parallel 15\pF$), que introducen una leve derivación capacitiva a $1\kHz$.
    \end{enumerate}
\end{itemize}
